\section{Experiments}
\label{sec:experiments}

\subsection{Dataset and Setup}

\noindent\textbf{Dataset.}
We evaluate on the MULTIAQUA dataset~\cite{muhovic2025multiaqua}, which provides synchronized RGB, thermal, and LiDAR data captured from unmanned surface vehicles.
The dataset contains four semantic classes: static obstacle, dynamic obstacle, water, and sky.
The training and validation sets consist of daytime images (145 validation images), while the test set contains only nighttime images evaluated through the MACVi challenge server~\cite{bovcon2024macvi}.
This setup requires models trained on daytime data to generalize to nighttime conditions where RGB quality degrades severely.

\noindent\textbf{Implementation details.}
We use SAM2~\cite{ravi2024sam2} with a Hiera-B+~\cite{ryali2023hiera} backbone as the foundation model, with LoRA~\cite{hu2022lora} rank $r{=}4$ and $N{=}3$ Soft-MoE experts.
All images are resized to $1024 \times 1024$.
We train with AdamW optimizer (learning rate $6 \times 10^{-4}$, weight decay $0.01$) using a warmup-polynomial schedule (10 warmup epochs, power $0.9$).
The batch size is 1 per GPU and we use Online Hard Example Mining (OHEM) cross-entropy loss with the prototype segmentation loss~\cite{chen2025memorysam} (\cref{sec:training}).
SAM2's image encoder is frozen; only the Soft-MoE LoRA layers, cross-modal fusion head, quality gating network, prompt encoder, and mask decoder are trained.
The quality gating loss weight is $\lambda_\text{gate} = 0.5$.
To bridge the day-to-night domain gap, we apply physics-inspired RGB augmentations during training~\cite{muhovic2025multiaqua}: brightness and gamma reduction simulating low-light conditions, complementary random masking~\cite{shin2024crm}, and stochastic full-channel zeroing, while leaving thermal and LiDAR unchanged.
Training is conducted with PyTorch DDP on NVIDIA GPUs.

\noindent\textbf{Evaluation metric.}
We report mean Intersection-over-Union (mIoU) averaged over the four semantic classes, on both the daytime validation set and the nighttime test set.


\subsection{Comparison with State of the Art}

\cref{tab:main_results} compares our method against prior approaches on the MULTIAQUA benchmark.
The baselines include CMNeXt~\cite{zhang2023delivering}, MMSFormer~\cite{reza2024mmsformer}, and StitchFusion~\cite{li2024stitchfusion} with various training strategies from~\cite{muhovic2025multiaqua}: suffix ``-D'' denotes double-pass training (processing each modality pair twice with swapped order), and ``-H'' denotes a multi-head decoder.
All baseline results are taken from~\cite{muhovic2025multiaqua}.

\begin{table}[t]
  \caption{Comparison on the MULTIAQUA benchmark. Val mIoU (daytime) and test mIoU (nighttime) are reported. ``D'' = double-pass training, ``H'' = multi-head decoder. Baseline results from~\cite{muhovic2025multiaqua}. Best results in \textbf{bold}, second best \underline{underlined}.}
  \label{tab:main_results}
  \centering
  \setlength{\tabcolsep}{3.5pt}
  \begin{tabular}{@{}lccc@{}}
    \toprule
    Method & Val mIoU & Test mIoU & Obs.\ (test) \\
    \midrule
    CMNeXt$_\text{RGB}$~\cite{zhang2023delivering}  & \textbf{94.31} & 38.23 & 11.87 \\
    CMNeXt~\cite{zhang2023delivering}                & 88.78 & 50.52 & 10.86 \\
    CMNeXt-H                                         & 89.97 & 50.05 & 11.10 \\
    CMNeXt-D                                         & 92.95 & 72.24 & 32.68 \\
    CMNeXt-DH                                        & 93.58 & \underline{74.25} & \underline{37.25} \\
    \midrule
    MMSFormer~\cite{reza2024mmsformer}               & 85.68 & 40.11 & 10.37 \\
    MMSFormer-H                                      & 88.36 & 45.70 & 9.37 \\
    MMSFormer-D                                      & 89.98 & 64.12 & 13.36 \\
    MMSFormer-DH                                     & 88.21 & 63.09 & 15.89 \\
    \midrule
    StitchFusion~\cite{li2024stitchfusion}           & 93.57 & 41.86 & 14.98 \\
    StitchFusion-H                                   & 91.14 & 42.95 & 19.73 \\
    StitchFusion-D                                   & 89.81 & 74.23 & 36.89 \\
    StitchFusion-DH                                  & 91.59 & 73.59 & 34.38 \\
    \midrule
    MemorySAM~\cite{chen2025memorysam}               & \TODO{XX.X} & \TODO{XX.X} & \TODO{XX.X} \\
    Ours                                             & \TODO{XX.X} & \TODO{XX.X} & \TODO{XX.X} \\
    \bottomrule
  \end{tabular}
\end{table}

% TODO: Fill in MemorySAM and Ours results when P24 training completes.
% P9 hardaug8_physaug ep131: Val 93.54 / Test 70.41 / Obstacle 32.85
% P24: pending CE-based retraining

\cref{tab:main_results} shows that among the baselines, the double-pass training strategy (``-D'') provides the largest improvement for nighttime generalization, with CMNeXt-DH achieving 74.25\% test mIoU.
% TODO: Add analysis of our results vs baselines when available.


\subsection{Ablation Studies}

\noindent\textbf{Effect of night augmentation.}
\cref{tab:ablation_nightaug} validates the importance of night-robust augmentation.
Without any RGB corruption, the model achieves high daytime validation mIoU but suffers catastrophic failure on the nighttime test set.
Night augmentation substantially bridges this gap by exposing the model to simulated degradation patterns during training.

\begin{table}[t]
  \caption{Ablation on night-robust RGB augmentation.}
  \label{tab:ablation_nightaug}
  \centering
  \begin{tabular}{@{}lcc@{}}
    \toprule
    Night Aug & Val mIoU & Test mIoU \\
    \midrule
    \xmark & 93.10 & 35.93 \\
    \cmark & \TODO{XX.X} & \TODO{XX.X} \\
    \bottomrule
  \end{tabular}
\end{table}

\noindent\textbf{Effect of quality-aware memory gating.}
To isolate the contribution of quality-aware memory gating, we compare our full model against the baseline without the gating module (equivalent to the MemorySAM backbone with our other components).

% TODO: Fill in after P24 training completes
\begin{table}[t]
  \caption{Ablation on quality-aware memory gating. ``QMG'' = quality-aware memory gating.}
  \label{tab:ablation_qmg}
  \centering
  \begin{tabular}{@{}lcc@{}}
    \toprule
    QMG & Val mIoU & Test mIoU \\
    \midrule
    \xmark (P9 baseline) & 93.54 & 70.41 \\
    \cmark (Ours)        & \TODO{XX.X} & \TODO{XX.X} \\
    \bottomrule
  \end{tabular}
\end{table}

% \subsection{Qualitative Results}
% TODO: Add visualization of segmentation results on daytime and nighttime samples.
% Show quality maps for each modality and how they differ across conditions.
